% ==================================================
% CUBIERTA
% Fondo oscuro y una familia de gaussianas que se
% estrechan hasta convertirse en una delta de Dirac.
% ==================================================
\begin{pagina}
\thispagestyle{empty}
\begin{tikzpicture}[remember picture,overlay]

  % Fondo
  \fill[darkcolor] (current page.south west) rectangle (current page.north east);

  % Franja lateral fina
  \fill[maincolor] ([xshift=1.6cm]current page.south west)
        rectangle ([xshift=1.72cm]current page.north west |- current page.north west);

  % Gaussianas -> delta
  \begin{scope}[shift={([yshift=-4.2cm]current page.center)}]
    \draw[maincolor!50!white, opacity=0.35, line width=0.5pt]
      (-9,0) -- (9,0);
    \foreach \sig/\opa in {1.60/0.22, 1.20/0.30, 0.90/0.40, 0.65/0.52,
                           0.45/0.66, 0.30/0.82, 0.20/1.00}{
      \draw[color=maincolor!70!white, opacity=\opa, line width=0.9pt,
            domain=-9:9, samples=240, smooth, variable=\x]
        plot ({\x},{1.1/\sig*exp(-(\x*\x)/(2*\sig*\sig))});
    }
    % Delta de Dirac
    \draw[white, line width=1.1pt, -{Stealth[length=7pt,width=6pt]}] (0,0) -- (0,6.6);
  \end{scope}

  % Cabecera
  \node[anchor=north west, text=maincolor!35!white, inner sep=0pt]
    at ([xshift=2.8cm,yshift=-3cm]current page.north west)
    {\scshape\large Mecánica teórica};

  % Título
  \node[anchor=north west, text=white, align=left, inner sep=0pt, text width=15cm]
    at ([xshift=2.8cm,yshift=-4.6cm]current page.north west)
    {{\fontsize{44}{50}\selectfont\spaceskip=0pt\relax\bfseries Mecánica de\\[2pt] Distribuciones}\\[1.3cm]
     {\color{maincolor!45!white}\rule{3cm}{1.4pt}}\\[0.6cm]
     {\fontsize{14}{20}\selectfont\spaceskip=0pt\relax\color{white!85!maincolor}
      Fundamentos matemáticos\\ y formulación hamiltoniana}};

  % Autor
  \node[anchor=south west, text=white, inner sep=0pt]
    at ([xshift=2.8cm,yshift=2.4cm]current page.south west)
    {\scshape\Large Mario Sultan Romero};

\end{tikzpicture}
\end{pagina}
